\section{Del escalamiento en tiempo de prueba al escalamiento en tiempo de entrenamiento}

Las clases anteriores mejoraron la tasa de éxito con parámetros fijos mediante muestreo, búsqueda y un verifier. Esta clase plantea una pregunta más: \textbf{¿se pueden convertir esas reasoning trajectory exitosas en una señal de entrenamiento, de modo que se modifique la distribución de la salida única del modelo?}

AIME es el benchmark motivador principal. Contiene problemas recientes de matemáticas de competencia, lo que reduce el problema de que los benchmark antiguos estén contaminados por los datos de entrenamiento, y exige más de las cadenas de razonamiento largas, el retroceso y la autoverificación.

\videofigure{figures/aime.jpg}{AIME 2024/2025 se usa para medir la capacidad de razonamiento matemático reciente.}{00:01:16--00:03:28}

El cómputo en tiempo de entrenamiento y en tiempo de prueba son dos ejes que se pueden escalar simultáneamente: el entrenamiento hace que las trayectorias correctas sean más frecuentes en la distribución del modelo; el escalamiento en tiempo de prueba realiza una búsqueda más profunda o más amplia a partir de esa distribución.

\videofigure{figures/train-test-scaling.jpg}{Aumentar el cómputo de train-time y de test-time mejora el pass@1 de AIME.}{00:05:48--00:07:20}

\begin{importantbox}{El objetivo del escalamiento en tiempo de entrenamiento es remodelar la distribución de probabilidad}
El repeated sampling aumenta principalmente la probabilidad de "obtener al menos una vez la respuesta correcta"; el RL, en cambio, busca elevar en conjunto la probabilidad de las trayectorias con recompensa alta, para que la respuesta correcta aparezca con más frecuencia y consistencia. El resultado ideal es obtener la misma precisión con menos muestreo en tiempo de prueba.
\end{importantbox}

\subsection{Por qué el reasoning RL falla con facilidad}

En un CoT largo, el action space es enorme, la reward es escasa y la longitud de las trayectorias varía mucho. El entrenamiento también enfrenta problemas sistémicos como el reward hacking, el colapso de entropía, la varianza del gradiente, la presión de memoria y el throughput del muestreo en línea. Por eso el curso enfatiza que la "magia" de escalar el RL suele provenir de los detalles de implementación, no de una única función objetivo.

\subsection{Resumen de la sección}

El train-time scaling convierte al verifier o a la reward del entorno en actualizaciones de parámetros. Puede mejorar la estabilidad de una sola respuesta y de la votación mayoritaria, pero también puede reducir la diversidad de exploración; entender esta tensión es el hilo conductor común de STaR, GRPO y DAPO.
